% ── Section 7: Limitations and future work ─────────────────────────────
\section{Limitations and Future Work}
\label{sec:limitations_future}

The previous section identified the main transferable pattern: interpolation is a strong floor,
external predictors alone are not sufficient, TS-ICL adds skill, and the hardest regimes are
the distributional extremes. This final section states where those conclusions stop and what should be tested next.
The main deliverable that this reported framework gives CEAZA
is a reproducible way to decide when a reconstruction is credible, when it is only useful as a candidate
trajectory, and where new physical features or models are needed.

\subsection{Validation support and gap-length limits}
\label{sec:limits_validation}

The benchmark validates reconstruction skill only where artificial gaps can be created from
observed eligible data. This is a strong design choice, because every scored reconstruction is
compared with withheld truth, but it also limitats the validation support useful for the imputation of real gaps.

Real gaps sit in fact outside this validation logic. The real-gap reconstructions, shown for example in
Figure~\ref{fig:chl_real_gap_multigap}, are useful practical outputs for future research activities that benefit from continuous records, but they are not evidence of
method skill because the target was not observed inside those intervals. This distinction is
especially important for long outages: the shared 256-day station-wide gap during the 2020 pandemic, and other very long
real gaps, are longer than the artificial gaps used for model ranking. Any reconstruction of
such intervals is extrapolation from the validated regime, not a validated result.

\subsection{Extremes remain the main unresolved regime}
\label{sec:limits_extremes}

The most important limitation is not the ranking on a limited artificial-gap support, but the behaviour of the methods in the regimes that matter most scientifically. For chlorophyll, this means high-\chl{} event gaps, where Figure~\ref{fig:chl_event_limitation} shows that every method has higher error: even the best method still smooths sharp blooms and underestimates event peaks.
For oxygen, the corresponding limitation is again on the empirical tails of the distribution. TS-ICL improves in the benchmark because it performs better in the middle of
the oxygen distribution, but Figure~\ref{fig:oxygen_tail_diagnostics} shows weaker behaviour in the low and high tails, especially when tail conditions persist. The oxygen result therefore gives a
clear example of why a pooled MAE improvement is not enough for full operational trust: the average can
improve while the scientifically sensitive strata remain unreliable. Future work should therefore make event and tail performance part of the objective, not only a
post-hoc diagnostic. For chlorophyll, this could mean event-aware validation and modelling, with
losses or selection criteria that explicitly penalize missed bloom peaks. For oxygen, it means
separating isolated tail excursions from sustained tail episodes and testing whether models can
avoid regression toward the centre of the distribution. 

\subsection{Physical interpretation and covariate limits}
\label{sec:limits_physical}

The covariate analyses are a useful product to identify reconstructive signal from physical forcing predictors. For chlorophyll, it's wind/upwelling forcing that establishes useful signal, while for oxygen it's current information that proves being more informative, although it appears redundant once combined with
the broader physical covariate bundle. These findings do not
prove that the corresponding physical processes are causal, but are anyway useful because they
show which information channels improve reconstruction and indicates this approach as a future pipeline/research tool.

This distinction between causality and reconstruction improvement matters in Tongoy Bay because the relevant processes occur across scales. The
available gridded products describe regional surface chlorophyll, SST, winds and currents, but the
in-situ sensors measure conditions at a nearshore buoy inside a bay. Short-lived blooms, retention,
local circulation and aquaculture-scale variability may be only indirectly represented in the external products. A
negative reconstruction result for a predictor family therefore cannot be read as ruling out the
underlying process.

A natural extension is to design physical features more directly around the unresolved regimes. For
chlorophyll, this means testing whether wind history, local retention, recent cooling or
spatial patchiness can better anticipate bloom formation and decay. For oxygen, it means building
features that distinguish sustained low- or high-oxygen episodes from isolated excursions, and
testing whether current variables contribute unique information once the broader physical
state is controlled.

\subsection{Model and uncertainty limits}
\label{sec:limits_models_uncertainty}

TS-ICL is promising here because it improves in the benchmark for both targets without station-specific fitting or fine-tuning, and because it can use target history and covariates as
time-indexed channels rather than relying on a large hand-built tabular feature matrix. But this
does not make it a universal solution, as only one pretrained time-series imputation foundation model was tested,
and the results do not establish that TS-ICL is the best architecture available for coastal
imputation, tabular foundation models being a first natural extension. Anyway, in this scenarios where datasets have a limited size and covariates on coarse grids may not explain well nearshore measurements, the results establish that this class of zero-shot models deserves comparison against classical baselines.

The uncertainty evidence is also incomplete. For chlorophyll for instance, TS-ICL satellite-proxy quantile
intervals are close to nominal overall, but event-day calibration is weaker. Future work should treat interval calibration as a first-class benchmark output.

A further model-level limitation is boundary behaviour. Boundary-anchored methods such as linear
interpolation and gap-edge residual models are constrained by observations at the edges of an
internal gap. TS-ICL uses a broader context window and covariate channels, and is not explicitly
forced to match the nearest observations at the gap boundary. This can be an advantage when local
endpoints are misleading, but it can also create edge discontinuities. A practical next step is to
test hybrid methods that preserve TS-ICL's interior reconstruction while blending more conservatively near the gap edges.

\subsection{Extensions of the benchmark workflow}
\label{sec:future_extensions}

The main reusable output of this report is the benchmark workflow. A new sensor or station can be
processed through the same sequence: define the daily target and eligibility rule, audit missingness,
construct artificial gaps, apply leakage-controlled comparators, stratify errors by scientifically
relevant regimes, and only then produce real-gap reconstructions if the evidence supports them.

Several extensions follow directly from this structure. The first is spatial transfer: applying the
same benchmark to other CEAZA stations would test whether the model hierarchy found at Tongoy is
specific to this site or network-wide. The second is target transfer: additional variables such as
temperature, salinity, turbidity, pH or other biogeochemical sensors could be evaluated with the
same design. The third is long-gap validation: longer artificial gaps are rare because they require
long uninterrupted observed runs, so pooling multiple stations or alternative validation designs
may be needed to evaluate outages beyond the current support.